\documentclass[11pt,letterpaper]{article}
\usepackage[T1]{fontenc}
\usepackage{lmodern}
\usepackage[margin=1in]{geometry}
\usepackage{amsmath,amssymb,amsthm,mathtools}
\usepackage{booktabs,array,microtype,enumitem}
\usepackage[colorlinks=true,linkcolor=black,citecolor=black,urlcolor=blue]{hyperref}
\usepackage{fancyhdr}
\hypersetup{pdftitle={A Finite Positive Amplitude for Stressed Kunz Words},pdfsubject={Research report 274},pdfauthor={},pdfcreator={pdfLaTeX}}
\pdfmapfile{+lm.map}\pdfmapfile{+cm.map}\pdfmapfile{+symbols.map}\pdfmapfile{+euler.map}
\pdfinfoomitdate=1\pdftrailerid{}\pdfsuppressptexinfo=15
\pagestyle{fancy}\fancyhf{}
\fancyhead[L]{\small Finite amplitude for stressed Kunz words}
\fancyhead[R]{\small Report 274}\fancyfoot[C]{\thepage}
\setlength{\headheight}{14pt}\setlength{\emergencystretch}{2em}
\setlist{itemsep=3pt,topsep=5pt}\numberwithin{equation}{section}
\newtheorem{theorem}{Theorem}[section]
\newtheorem{lemma}[theorem]{Lemma}
\newtheorem{proposition}[theorem]{Proposition}
\newtheorem{corollary}[theorem]{Corollary}
\theoremstyle{definition}
\newtheorem{definition}[theorem]{Definition}
\newtheorem{remark}[theorem]{Remark}
\newcommand{\N}{\mathbb N}\newcommand{\Z}{\mathbb Z}
\newcommand{\eps}{\varepsilon}\newcommand{\wt}{\operatorname{wt}}
\newcommand{\dist}{\operatorname{dist}}
\newcommand{\cS}{\mathcal S}\newcommand{\cU}{\mathcal U}
\newcommand{\cV}{\mathcal V}\newcommand{\cT}{\mathcal T}
\newcommand{\Hbin}{H_{\mathrm{bin}}}
% Added in the write (7 October 2026).
\newcommand{\file}[1]{\nolinkurl{#1}}
\newenvironment{writenote}[1][]{\par\smallskip\begingroup\small
 \noindent\textbf{[write]}\if\relax\detokenize{#1}\relax\else~\textbf{#1.}\fi\ }%
 {\par\endgroup\smallskip}
\title{\textbf{A Finite Positive Amplitude\\for Stressed Kunz Words}\\[5pt]
\large A Self Contained Proof with an Exponential Remainder}
\author{Research report 274}\date{6 October 2026}
\begin{document}\maketitle
\begin{abstract}
Let $s_g$ count words on $\{1,2,3\}$ ending in $3$ and satisfying the Kunz inequalities, with genus equal to the sum of their letters. We prove
\[
s_g=C_s\rho^g+O(\rho_2^g),\qquad 0<C_s<\infty,\qquad 0<\rho_2<\rho,
\]
where $\rho^{-1}=\xi$ is the positive root of
$(\xi^2+\xi^3)(\xi+\xi^2+\xi^3)=1$.
An exact boundary renewal gives a convergent positive series for
$C_s=N(\xi)/(\xi P'(\xi))$. Its convergence follows from an exact activity transformation and an elementary weighted container and matching estimate. The complementary words, having a $1$ at or before one third of their length, have exponentially small total critical weight. The decisive additional loss comes from translated forbidden pairs and a deterministic selection of disjoint matching blocks. All combinatorial estimates, preimage losses, and analytic steps are proved here; earlier root and ratio theorems are not proof inputs. The result establishes the stressed part of Zhu's polynomial-prefactor formulation with exponent zero, and gives a two-pole expansion for the depth-at-most-three count. We also obtain a shrinking-error inverse threshold and the vanishing of every fixed inverse-power correction. No finite amplitude for depth at least four, numerical value of $C_s$, practical onset, or full subdominant expansion is asserted. An original bounded exact-arithmetic companion checks the finite identities and constructions.
\end{abstract}
\begin{writenote}[Status and trust boundary, added 7 October 2026]
This is Report~274 of an external AI-assisted research session, filed in
ProveIt as the research report \file{numerical-semigroup-stressed-amplitude}
of the enumerative-combinatorics collection (provenance, the sources as the
write read them, what was checked, relation to the repository, the collected
non-claims and the reading conventions are in
Section~\ref{ska:sec:provenance}). It is unrefereed and not formalized. Its
author line reads ``Research report 274'' and its PDF author field is empty:
it names no person or tool. \emph{The main theorem is not independently
verified.} Theorem~\ref{ska:thm:main} rests on a long self-contained argument
(the exact renewal, the activity transformation, a loop-aware container
lemma, dense-window standardization and the early-one penalty of
Section~\ref{ska:sec:early}); the write read every step and found no error,
and every exact finite identity it tested holds, but a reading is not a
verification, and the source itself asks for an independent check of the
all-length argument. The constants of the proof are existential and
astronomically weak (the early-one saving is below $10^{-4}$ per letter),
so no finite computation tests the conclusion: the finite data are still far
from the asserted limit (notes in Sections~\ref{ska:sec:theorem}
and~\ref{ska:sec:questions}). If the theorem holds, it proves the $s_g$ parts
of Zhu's Conjecture~7.3 with exponent $\alpha'=0$, against the numerical
expectation stated there. The write adds Remarks~\ref{ska:rem:transseries}
and~\ref{ska:rem:oeis} and dated notes; no statement, proof or number of the
source was changed.
\end{writenote}
\clearpage
\setcounter{tocdepth}{1}
\makeatletter\renewcommand{\l@section}{\@dottedtocline{1}{0em}{2.2em}}\makeatother
\tableofcontents\clearpage

\section{The theorem and its scope}\label{ska:sec:theorem}
A \emph{stressed word} is a nonempty word $w=w_1\cdots w_L$ on $\{1,2,3\}$, ending in $3$, such that
\begin{equation}\label{ska:eq:unwrapped}
w_i+w_j\ge w_{i+j}\qquad(1\le i,j,\ i+j\le L).
\end{equation}
Repeated source indices are permitted. On this alphabet, the wrapped Kunz inequalities are automatic, and the only forbidden unwrapped pattern is two $1$'s whose index sum carries $3$. Set
\[
g(w)=\sum_{i=1}^Lw_i,\qquad
s_g=\#\{w\text{ stressed}:g(w)=g\},\qquad
\cS(z)=\sum_{g\ge0}s_gz^g.
\]
Thus $s_0=s_1=s_2=0$. Define
\begin{align}
A(z)&=z^2+z^3,& B(z)&=z+z^2+z^3,& C(z)&=z+z^2,\label{ska:eq:ABC}\\
P(z)&=A(z)B(z)=z^3+2z^4+2z^5+z^6,
&P(\xi)&=1,& \rho&=\xi^{-1}.\label{ska:eq:P}
\end{align}
The positive solution $\xi$ is unique because $P$ is strictly increasing on the positive real axis, starts at zero, and is unbounded. We shall use the exact inequalities
\begin{equation}\label{ska:eq:rootbracket}
\varphi^{-1}<\frac58<\xi<\frac{33}{50}<\frac23<1,
\qquad \varphi=\frac{1+\sqrt5}{2}.
\end{equation}
Indeed, $5/8+(5/8)^2>1$ gives the first comparison, direct substitution gives $P(5/8)<1$, and
\begin{equation}\label{ska:eq:rootcert}
P(33/50)-1=\frac{1737269}{15625000000}>0.
\end{equation}
Equivalently, $\rho>1$ solves $\rho^6-\rho^3-2\rho^2-2\rho-1=0$.

For $r\ge1$, and a set $J\subseteq\{0,\ldots,r\}$, let
\[
J+J=\{j+j':j,j'\in J\},\qquad
u(J)=|(J+J)\cap\{0,\ldots,r-1\}|.
\]
The sumset includes sums of an element with itself. Define the nonnegative polynomials
\begin{equation}\label{ska:eq:Dr}
D_r(z)=A(z)^r
\sum_{\substack{J\subseteq\{0,\ldots,r\}\\0\in J,\ r\notin J+J}}
z^{|J|+3}A(z)^{r+1-|J|}C(z)^{u(J)}B(z)^{r-u(J)},
\end{equation}
and the formal series
\begin{equation}\label{ska:eq:N}
N(z)=z^3(1+A(z))+\sum_{r\ge1}D_r(z).
\end{equation}

\begin{theorem}[Finite positive stressed amplitude]\label{ska:thm:main}
The series in~\eqref{ska:eq:N} converges absolutely and uniformly on $|z|\le33/50$. There exist constants $K<\infty$ and $0<\rho_2<\rho$ such that, for all integers $g\ge0$,
\begin{equation}\label{ska:eq:main}
\left|s_g-C_s\rho^g\right|\le K\rho_2^g,
\qquad
C_s=\frac{N(\xi)}{\xi P'(\xi)}\in(0,\infty).
\end{equation}
More precisely, for some $R_s>\xi$,
\begin{equation}\label{ska:eq:simplepole}
\cS(z)=\frac{C_s}{1-z/\xi}+H(z),
\end{equation}
where $H$ is analytic on $|z|<R_s$.
\end{theorem}

All constants in this article are existence constants unless an explicit value is displayed. In particular, the theorem does not provide a useful numerical approximation of $C_s$ or a practical genus from which its error term is small.

\subsection{How the proof is organized}
Partition the words into $U$, with no $1$ at or before $L/3$, and its complement $V$. Their generating functions satisfy the exact identity
\begin{equation}\label{ska:eq:proofmap}
\cS(z)=\underbrace{\frac{N(z)}{1-P(z)}}_{\cU(z)}+\cV(z).
\end{equation}
There are two separate convergence problems. First, the boundary polynomials $D_r$ must be summable beyond $\xi$. Their exact transformation into weighted stressed-word polynomials permits a strict subcritical estimate at $33/50$. Second, $V$ must have a radius of convergence greater than $\xi$. The ordinary critical matching estimate has a neutral factor, so a new loss is needed. An original early $1$ supplies linearly many necessary translated exclusions; a disjoint-block selection converts them into a strict exponential penalty. Only after both steps does the analytic pole extraction apply.

\subsection{Relation to the earlier literature}
The stressed-word encoding, the rational lower family with denominator $1-P$, and the Fibonacci filter for depth at most three have antecedents in Bacher and Zhu~\cite{Bacher,Zhu}. Bacher's discussion of central and terminal transition regions motivates separating a finite boundary from the two bulk blocks. The elementary fingerprint method used here belongs to the classical graph-container framework; see Samotij~\cite{Samotij}. Zhu introduced the dense-window standardization~\cite[Sections~5.4--5.5]{Zhu}; we use matching-preserving refinements and terminal-matching optimization from Reports~270 and~271~\cite{Report270,Report271}. All needed ingredients are proved in full below. No root-rate or consecutive-ratio conclusion from those reports is assumed, and no weighted graph-homomorphism theorem is needed for this three-letter argument.

Zhu's Section~7.3(c) asks for exponents $\alpha,\alpha'$ in polynomial-prefactor statements for the depth-at-least-four count and for $s_g$. Its formal statement does not require the exponents to be positive. Theorem~\ref{ska:thm:main} establishes the stressed part with $\alpha'=0$. This differs from the numerical suggestion there that the exponents lie between $1.6$ and $1.7$; it does not contradict a proved theorem or the formal existence assertion. The depth-at-least-four exponent $\alpha$, its amplitude, and its ratio to $s_g$ remain open here.

The present proof is self contained apart from elementary complex analysis. No comprehensive historical-priority claim is made. In view of the difference from the finite-data extrapolation, independent external checking of the all-length argument would be especially valuable.
\begin{writenote}[Zhu's conjecture and his finite data, 7 October 2026]
The write read Zhu's arXiv version~3 (7 September 2023, 20 pages). The
passage cited as Section~7.3(c) is Conjecture~7.3(c) (p.~19), stated jointly
for the depth-at-least-four count $\hat n_g=h_g$ and for $s_g$: there are
constants $\alpha,\alpha'$, possibly equal, with $\hat n_g=g^{\alpha+o(1)}\rho^g$
and $s_g=g^{\alpha'+o(1)}\rho^g$; Zhu's $r_{1.51}$ is this report's $\rho$
(the positive root of $x^6-x^3-2x^2-2x-1$, defined in his Section~2). The
surrounding text says that, on the numerical evidence, exponential
asymptotics for $s_g$ and $\hat n_g$ ``appears unlikely'' and that the
exponents seem to lie between $1.6$ and $1.7$. Theorem~\ref{ska:thm:main}
would prove the $s_g$ halves of parts (a)--(c) with $\alpha'=0$
(Corollary~\ref{ska:cor:stress} states the ratio form of part (b)); the
$\hat n_g$ halves stay open, as the source says. Zhu's Table~2 of $s_g$ for
$g\le95$ (p.~17), which the write's
own enumeration reproduces for $g\le26$, gives $s_g/\rho^g=0.3276$, $0.4937$,
$0.9137$, $1.4309$, $2.1860$ at $g=20$, $30$, $50$, $70$, $95$, still
increasing, with local exponent $d\log(s_g\rho^{-g})/d\log g$ of $1.16$,
$1.31$, $1.37$, $1.40$ at $g=40$, $60$, $80$, $95$. These data do not
contradict the theorem, whose remainder carries no effective constant, but
they show that its limit regime, if it is reached, lies beyond all available
counts.
\end{writenote}

\subsection{Provenance, sources and reading conventions}\label{ska:sec:provenance}
\begin{writenote}[Added 7 October 2026]\raggedright
\emph{Provenance.} The archive \file{Report274_Stressed_Kunz_Finite_Amplitude.zip}
($448{,}517$ bytes, 9 files, wrapper directory \file{Report274/}) arrived
alone in \file{764740f07} (6~October 2026); it is not part of the session
bundle, and it was placed in \file{2680aae95} in the enumerative-combinatorics
collection, since no OEIS sequence is its object. The text is the delivered
\file{article.tex} (824 lines, 21 pages, dated 6~October 2026), shipped under
the same name. The programs are shipped under \file{code/}: the builder
\file{build.py} as \file{code/build.py}, \file{companion/exact_checks.py} as
\file{code/companion-exact_checks.py} and the two test files as
\file{code/tests-test_build.py} and \file{code/tests-test_companion.py};
\file{SOURCES.md} as delivered. Not shipped, and retrievable from the arrival
commit: the delivered \file{article.pdf} (replaced by a build of this text),
the checksum manifest \file{MANIFEST.sha256} (8 entries, all verified at the
write) and the delivered \file{README.md}, which the report README replaces.
The package records no ProveIt commit; it carries no ``prepared for private
review'' line and no personal data. It is a single-source report. It names
Reports~270 and~271 as earlier reports of the same research sequence; neither
was delivered, and if they arrive they continue this report. The write
prefixed the 96 delivered labels with \file{ska:}, a prefix used
nowhere else in the repository, and updated the references to them
(51 \verb|\eqref|, 13 \verb|\ref|); it added this subsection, the notes marked [write] and
Remarks~\ref{ska:rem:transseries} and~\ref{ska:rem:oeis}, each the last
statement of its section, and inserted no numbered display, so every
section, statement and equation of the source keeps its delivered number.

\emph{The sources, as the write read them} (7~October 2026): Zhu's
arXiv:2202.05755v3 (note at the end of Section~1.2; Definition~3.7 of a
stressed word, Section~4.1 with the lower family $x^3(x+1)(x^2+x+1)$,
Conjecture~7.3 and Table~2 of $s_g$, $t_g$, $\hat n_g$); OEIS A007323 in the
internal format, and OEIS searches for $s_g$ and $t_g$
(Remark~\ref{ska:rem:oeis}); the transseries volume
\path{Analysis/Transseries/docs/series-and-transseries/Transseries_And_Inversion/transseries_and_inversion.tex}
(\file{p0:def:model}, \file{p0:thm:lambert-core}, \file{p0:def:three-inverses},
\file{p0:thm:staircase}). Not read by the write: Bacher, Samotij, Zhu's
Sections~5.4--5.5 in detail, and the undelivered Reports~270 and~271, which
the source says are not proof dependencies.

\emph{What was checked at intake.} The intake dossier read the manuscript,
found no demonstrably wrong claim, ran the companion's exact checks and its
13 regression tests on a copy (they pass; the Linux-only builder and its
tests were not run), and recorded that the counts to $g=36$ and the boundary
series to $r=23$ are pre-asymptotic. \emph{What the write checked.} It read
every proof and found no error; in particular it checked the block
decomposition of Proposition~\ref{ska:prop:renewal}, the sequential bookkeeping
in the proof of Proposition~\ref{ska:prop:transform} (offset~$0$ carries the
endpoint multiplier $u\gamma$), the fingerprint reconstruction and the
incidence count in Lemma~\ref{ska:lem:containers}, the independence of the
standardized prefix one-set, the eligibility optimization of
Lemma~\ref{ska:lem:product}, the three cases of
Theorem~\ref{ska:thm:early} (including $a-|H|>L/16-1$, the source count
$|H|-q\ge L/12$, the degree bound four and the selection of at least $m/8$
block-disjoint relations), and the residue computations of
Sections~\ref{ska:sec:analytic} and~\ref{ska:sec:twopole}. With its own code
(Python 3.14.4, exact integers and fractions; SymPy 1.14.0; mpmath 1.3.0), not
importing the delivered programs, it
\begin{itemize}\raggedright
\item enumerated all Kunz words on $\{1,2,3\}$ of genus $g\le26$: $s_g$ and
$t_g$ equal Zhu's Table~2, the filter~\eqref{ska:eq:filter} holds (and holds
in that table for all $3\le g\le95$), and the separated class $U$ equals
$N/(1-P)$ coefficientwise, also with the length refinement
of~\eqref{ska:eq:refinedrenewal} (at $g=26$: $10{,}146$ separated and
$10{,}806$ early-one words of $s_{26}=20{,}952$);
\item verified~\eqref{ska:eq:transform} as a polynomial identity for $r\le6$
and modulo $z^{27}$ for $r\le8$, the general identity~\eqref{ska:eq:generalrenewal}
at three random rational triples for $r\le8$, and~\eqref{ska:eq:iterates} for
$n\le3$;
\item checked exactly every rational certificate: \eqref{ska:eq:rootbracket}
($P(5/8)=209625/262144$), \eqref{ska:eq:rootcert},
\eqref{ska:eq:actcert1}--\eqref{ska:eq:triplecert}, the formula for $E(x)^2$,
its derivative~\eqref{ska:eq:Ederivative}, $E(2/3)^2=1120000/1121931$,
$T_0(2/3)=64/81$, and the identity behind $L_2>L_b$;
\item computed $\xi=0.6599822256094898538\ldots$, $\rho=1.5151923206363711598\ldots$
(Zhu's $r_{1.51}$), $\xi P'(\xi)=4.1282412080\ldots$, the critical factors
$\beta\gamma=0.7921\ldots$, $T_0=0.7646\ldots$, $E=0.9627\ldots$,
$L_1=1.9127\ldots$, $\eta=0.009838\ldots$, and the three components
$0.000790$, $0.0405$, $0.0000768$ of the saving~$\tau$
in~\eqref{ska:eq:tau}, which is therefore about $7.7\times10^{-5}$;
\item evaluated $D_r(\xi)=A(\xi)S_r(\cdot)$ for $r\le33$ by enumerating the
one-sets of the transformed words (note in Section~\ref{ska:sec:questions}).
\end{itemize}
It reran, on a fresh copy of the delivered layout,
\file{companion/exact_checks.py} and \file{tests/test_companion.py} with the
delivered flags (\texttt{-I -B -X int\_max\_str\_digits=640}, normal and
\texttt{-O}): the two companion outputs are identical, every count in the
delivered README's coverage list is reproduced, and the 13 tests pass in both
modes. The builder and its 33 tests need Linux descriptor features and were
not run. Floating computations are observations, not certificates.

\emph{Relation to the repository.} Before this placement the repository had
no report on numerical-semigroup enumeration by genus.
\file{numerical-semigroup-leaf-types} treats the same objects with a
structural question (leaves of large type in the semigroup tree) and
\file{a069762-pyramidal-frobenius} treats Frobenius numbers; no result is
shared and no reciprocal note is proposed. No Lean or Rocq file treats these
objects; no statement of this report is formalized, and its place in the
collection confers no formal status. The inverse is compared with the
transseries volume in Remark~\ref{ska:rem:transseries}.
\end{writenote}
\begin{writenote}[What is not claimed, collected from the source]
The stressed-word encoding, the rational lower family and the Fibonacci
filter are Bacher's and Zhu's; the dense-window standardization is Zhu's; the
container method is classical (Samotij). No comprehensive priority claim, no
external peer review. Not claimed: a numerical value of $C_s$ or a useful
approximation of it; a practical genus from which the error is small; any
amplitude, ratio or exponent for depth at least four; the sign or
noncancellation of the full secondary residual; the next singularity,
optimal remainder bases or a full subdominant expansion; a practical
numerical inverse; that finite tests prove the weighted criterion or the
three-case argument. \emph{The write adds:} its checks are exact or floating
computations; its transseries remark claims no novelty for any inversion.
\end{writenote}
\medskip
\begingroup\small
\noindent\emph{Reading conventions} (letters the source uses in more than one
sense, with the tempting false readings; no symbol was renamed).\par\smallskip
\renewcommand{\arraystretch}{1.1}
\noindent\begin{tabular}{@{}>{\raggedright\arraybackslash}p{0.95in}>{\raggedright\arraybackslash}p{5.3in}@{}}
\toprule
Symbol & Meanings\\
\midrule
$\alpha$, $\beta$, $\gamma$ & In Section~\ref{ska:sec:transform}: $\alpha=v+w$, $\beta=u+v+w$, $\gamma=u+v$; in Sections~\ref{ska:sec:standardization}--\ref{ska:sec:early}: $\beta=v+w$, $\gamma=u+v$, $W=u+v+w$ (the source flags this). \emph{False reading:} the $\beta$ of Section~\ref{ska:sec:transform} is the $W$ of Sections~\ref{ska:sec:standardization}--\ref{ska:sec:early}.\\
$A$, $B$, $C$, $P$, $\Pi$ & $A,B,C,P$ the fixed polynomials~\eqref{ska:eq:ABC}--\eqref{ska:eq:P}; $\Pi=\beta(\beta+u)$ the local pair factor, equal to $P(\xi)=1$ at the critical activities. $C_s$ the amplitude; $C$ also constants.\\
$D$, $N$, $n$ & $D_r$ the boundary polynomials, $N(z)$ their numerator series; $n=\lceil L/k\rceil$ the window length; $n_g$ all semigroups of genus $g$.\\
$H$, $R$ & $H$ a container, and $H(z)$ the analytic remainder~\eqref{ska:eq:analyticremainder}; $R$ the residual set of Lemma~\ref{ska:lem:containers}, $R_s$, $R'$ radii.\\
$L$, $M$ & $L$ the word length; $L_1,L_2,L_b$ eligibility gains; $M$ the preimage ratio~\eqref{ska:eq:M}, $M_0$ a constant in Section~\ref{ska:sec:inverse}.\\
$q$, $Q$ & $q$ the marked early one in Section~\ref{ska:sec:early}, and $q=\gamma/\beta$ in the proof of Proposition~\ref{ska:prop:transform}; $Q$ the window of $k$ three-positions.\\
$S$, $\cS$, $T$ & $S_r$ the weighted stressed polynomials~\eqref{ska:eq:weightedpoly}; $\cS(z)$ the generating function; $S_{\mathrm{lead}}$ the Fibonacci constant; $T_0$ the tail-pair factor, $\mathfrak T$ the activity map, $\cT(z)$ the depth-at-most-three series. \emph{False reading:} $S_r$ is not a coefficient of $\cS$.\\
$t$, $u$, $v$ & $t=\lfloor\sqrt k\rfloor$ in Lemma~\ref{ska:lem:containers} and $t=u/\alpha$ in Section~\ref{ska:sec:transform}; $u,v,w$ activities, $u(J)$ the sumset count in~\eqref{ska:eq:Dr} and $u_L$ the overwrite count~\eqref{ska:eq:overwritecount}; $v$ also an image word.\\
$\delta$, $\lambda$, $\rho$ & $\delta=1/48$ in Section~\ref{ska:sec:early}, $\delta_k$ the container slack, $\delta_*$ the inverse exponent; $\lambda=\rho_2/\rho$; $\rho=1/\xi$, $\rho_2$, $\rho_3$ remainder bases; $r$ the boundary parameter and $r=\varphi^{-1}$ in the proof of Corollary~\ref{ska:cor:twopole}.\\
\bottomrule
\end{tabular}
\endgroup

\section{Kunz words and the exact Fibonacci filter}\label{ska:sec:kunz}
For completeness, we connect the word count to numerical semigroups without importing an asymptotic theorem. A numerical semigroup is an additive submonoid $\Lambda\subseteq\N_0$ with finite complement. Its genus is $|\N_0\setminus\Lambda|$, its multiplicity is the least positive element $m$, and its conductor is the least $c$ such that $c+\N_0\subseteq\Lambda$. The depth is $\lceil c/m\rceil$; $\N_0$ has depth zero.

For each $1\le i<m$, let $k_i m+i$ be the least element of $\Lambda$ congruent to $i$ modulo $m$. Then $k_i\ge1$, and closure is equivalent to
\begin{align}
k_i+k_j&\ge k_{i+j} &&(i+j<m),\label{ska:eq:kunz1}\\
k_i+k_j+1&\ge k_{i+j-m} &&(i+j>m),\label{ska:eq:kunz2}
\end{align}
where $k_0=0$ and the case $i+j=m$ is automatic. Conversely, positive integers satisfying these inequalities define a semigroup by taking all multiples of $m$ and the progressions $k_i m+i+m\N_0$. The displayed inequalities check closure between every two residue classes; each progression has a least element, so the complement is finite. No positive integer below $m$ is included. This proves the bijection.

There are exactly $k_i$ missing nonnegative integers in residue class $i$, so the genus is $\sum_i k_i$. If $q=\max_i k_i$ and $j$ is the last index attaining $q$, the largest missing integer is $(q-1)m+j$, hence the depth is $q$. Thus depth at most three corresponds precisely to words on $\{1,2,3\}$ satisfying the Kunz inequalities. For these letters, the left side of~\eqref{ska:eq:kunz2} is at least three, so only~\eqref{ska:eq:kunz1} can restrict the word. Its only possible violation is $(1,1,3)$. This is the stressed condition when the final letter is three.

Let $t_g$ count semigroups of depth at most three, including $t_0=1$, and put $\cT(z)=\sum_{g\ge0}t_gz^g$. Let $n_g$ count all semigroups of genus $g$, and let $h_g=n_g-t_g$ count those of depth at least four.

\begin{proposition}[Classical Fibonacci filter]\label{ska:prop:filter}
As formal power series,
\begin{equation}\label{ska:eq:filter}
\cT(z)=\frac{1+\cS(z)}{1-z-z^2}.
\end{equation}
In particular, $s_g=t_g-t_{g-1}-t_{g-2}$ for $g\ge3$.
\end{proposition}
\begin{proof}
Every word with letters in $\{1,2\}$ is Kunz. A word containing a three is uniquely a stressed prefix, ending at its last three, followed by an arbitrary word on $\{1,2\}$. Conversely, appending such a tail to a stressed prefix preserves the inequalities: an unwrapped target in the tail is at most two, and every wrapped inequality on this alphabet is automatic. This is a bijection preserving genus. The arbitrary tail contributes $1/(1-z-z^2)$, and the empty stressed prefix accounts for the numerator $1$.
\end{proof}

\section{Exact renewal for the separated boundary}\label{ska:sec:renewal}
Let $U$ consist of stressed words of length $L$ whose first one, if present, occurs at $a>L/3$. For a second formal variable $y$, recording length, write
\[
\cU(z,y)=\sum_{w\in U}z^{g(w)}y^{|w|},\qquad \cU(z)=\cU(z,1).
\]

\begin{proposition}[Length-refined renewal]\label{ska:prop:renewal}
There is an exact identity of formal series
\begin{equation}\label{ska:eq:refinedrenewal}
\cU(z,y)=
\frac{yz^3(1+yA(z))+\displaystyle\sum_{r\ge1}y^{3r+2}D_r(z)}
{1-y^2P(z)}.
\end{equation}
In particular, $\cU(z)=N(z)/(1-P(z))$.
\end{proposition}
\begin{proof}
First count the words whose ones are all strictly after the midpoint. For odd length $L=2m+1$, the first $m$ positions have alphabet $\{2,3\}$, the next $m$ positions are arbitrary, and the terminal position is three. Their weight is $z^3P(z)^m$. For even length $L=2m+2$, there are $m+1$ initial positions on $\{2,3\}$ and $m$ subsequent arbitrary positions before the terminal three, giving $z^3A(z)P(z)^m$. No sum of two one-positions reaches a position in the word. These disjoint odd and even families, including words with no ones, give the first numerator in~\eqref{ska:eq:refinedrenewal}.

Every remaining word in $U$ has $L/3<a\le L/2$. Put $r=L-2a$. The value $r=0$ is impossible because the one at $a$ would add to itself to reach the terminal three. Therefore $r\ge1$, and $a>L/3$ is equivalent to $a>r$. The consecutive blocks are
\begin{center}
\begin{tabular}{@{}lll@{}}
\toprule
Positions & Size & Role\\
\midrule
$1,\ldots,a-1$ & $a-1$ & arbitrary $2/3$ letters\\
$a,\ldots,a+r$ & $r+1$ & active block starting with $1$\\
$a+r+1,\ldots,2a-1$ & $a-r-1$ & arbitrary letters\\
$2a,\ldots,2a+r=L$ & $r+1$ & target block ending with $3$\\
\bottomrule
\end{tabular}
\end{center}
Only ones in the active block can participate in a sum at most $L$: a later one, added to the smallest one-position $a$, already gives a sum greater than $L$. In particular, $3a>L$ excludes target-block ones as sources.

Let $J$ be the offsets of the ones in the active block. Then $0\in J$. Every target offset in $J+J$ is forbidden to carry three, and every other nonterminal target is arbitrary. The terminal three requires $r\notin J+J$, which also forces $r\notin J$. For fixed $J$, the active and target blocks therefore have weight
\[
z^{|J|+3}A^{r+1-|J|}C^{u(J)}B^{r-u(J)}.
\]
The other two blocks have weight $A^{a-1}B^{a-r-1}$. Writing $a=r+1+j$ gives $L=3r+2+2j$ and
\[
\sum_{a\ge r+1}y^{2a+r}A^{a-1}B^{a-r-1}
=\frac{y^{3r+2}A^r}{1-y^2AB}.
\]
This yields the second numerator in~\eqref{ska:eq:refinedrenewal}. The first-one location, $r$, and the active one-set are uniquely determined by the word. There are no marked-word multiplicities in this identity.
\end{proof}

The minimal length attached to $D_r$ is $3r+2$, so the bivariate formal sum is locally finite. The univariate genus sum is also locally finite because every letter contributes at least one. These observations justify the identities before any convergence is known.

\section{The exact transformation of activities}\label{ska:sec:transform}
For positive activities $(u,v,w)$, let
\begin{equation}\label{ska:eq:weightedpoly}
S_r(u,v,w)=\sum_{\substack{a_1\cdots a_r\text{ stressed}}}
u^{\#\{i:a_i=1\}}v^{\#\{i:a_i=2\}}w^{\#\{i:a_i=3\}}.
\end{equation}
This is a homogeneous polynomial of degree $r$ with nonnegative integer coefficients. The same notation will also be used for its formal polynomial evaluation.

\begin{proposition}[Boundary activity identity]\label{ska:prop:transform}
For every $r\ge1$ there is a polynomial identity
\begin{equation}\label{ska:eq:transform}
D_r(z)=A(z)S_r\bigl(A(z)^2,A(z)^2C(z),A(z)^2z^3\bigr).
\end{equation}
More generally, put $\alpha=v+w$, $\beta=u+v+w$, and $\gamma=u+v$, and define
\begin{equation}\label{ska:eq:Tmap}
\mathfrak T(u,v,w)=(u\alpha\gamma,\alpha^2\gamma,\alpha^2w).
\end{equation}
If $\cU(u,v,w)$ is the weighted sum over the separated-boundary class, then formally
\begin{equation}\label{ska:eq:generalrenewal}
\cU(u,v,w)=
\frac{w(1+\alpha)+u\gamma\displaystyle\sum_{r\ge1}S_r(\mathfrak T(u,v,w))}
{1-\alpha\beta}.
\end{equation}
\end{proposition}
\begin{proof}
We first verify the general boundary contribution. In the proof of Proposition~\ref{ska:prop:renewal}, replace the weights $z,z^2,z^3$ by $u,v,w$. The counterpart of $D_r$ is
\[
\Delta_r=\alpha^r
\sum_{\substack{0\in J\subseteq\{0,\ldots,r\}\\r\notin J+J}}
u^{|J|}w\alpha^{r+1-|J|}\gamma^{u(J)}\beta^{r-u(J)}.
\]
For the algebraic verification, use temporarily the positive quantities
\[
b=\alpha^2\beta,\qquad t=u/\alpha,\qquad q=\gamma/\beta.
\]
Then
\begin{equation}\label{ska:eq:sequentialsum}
\Delta_r=w\alpha b^r\sum_J t^{|J|}q^{u(J)}.
\end{equation}
The position zero, necessarily selected and covered, contributes $tq$. Write $J'=J\setminus\{0\}$ and process $i=1,\ldots,r-1$ in order. Whether $i\in J'+J'$ is determined by earlier positive selected positions. If $i$ is selected, its contribution after attaching one factor $b$ is $btq=u\alpha\gamma$. If it is not selected but is already covered by $J'+J'$, its contribution is $bq=\alpha^2\gamma$. If it is neither selected nor covered, its contribution is
\[
b=\alpha^2\gamma+\alpha^2w.
\]
These are exactly the choices of letter $1$, letter $2$ at a covered position, and letters $2$ or $3$ at an uncovered position in a stressed word with activities~\eqref{ska:eq:Tmap}. Selecting a covered position is permitted: a one at the target of two earlier ones never violates a Kunz inequality.

The final position must be three, of transformed weight $\alpha^2w$, and this is allowed precisely when $r\notin J'+J'$. The original condition $r\notin J+J$ additionally forbids selecting $r$; the terminal-three condition does the same in the transformed word. The endpoint multiplier left over from~\eqref{ska:eq:sequentialsum} is
\[
\frac{w\alpha\,b\,tq}{\alpha^2w}=u\gamma.
\]
Consequently $\Delta_r=u\gamma S_r(\mathfrak T(u,v,w))$. Although the calculation used positive quotients, both sides are polynomials, so equality on all positive triples proves the polynomial identity. Substitution into the exact renewal proves~\eqref{ska:eq:generalrenewal}.

For $(u,v,w)=(z,z^2,z^3)$, we have $u\gamma=zC(z)=A(z)$ and $u\alpha\gamma=A(z)^2$, giving~\eqref{ska:eq:transform}. The first endpoint check is
\[
D_1(z)=z^4A(z)^2C(z)=z^3A(z)^3
=A(z)S_1\bigl(A(z)^2,A(z)^2C(z),A(z)^2z^3\bigr).
\]
\end{proof}

The transformation is an exact identity for the separated class, not a renewal identity for all stressed words. One useful algebraic feature is that $\gamma/\alpha$ is invariant under $\mathfrak T$, since both $\alpha$ and $\gamma$ are multiplied by $\alpha\beta$. For example,
\begin{align}
\mathfrak T^n(z,z^2,z^3)
&=\omega_n\bigl(1,z+z^2+\cdots+z^{n+1},z^{n+2}\bigr),\label{ska:eq:iterates}\\
\omega_0&=z,\qquad
\omega_{n+1}=z\omega_n^3\left(\frac{1-z^{n+2}}{1-z}\right)^2.\nonumber
\end{align}
These formulas follow by substitution, with the quotient interpreted as a finite geometric sum. They alone do not establish convergence of the full weighted stressed series. The next estimates supply that essential analytic input.

\section{An elementary container estimate with loops}\label{ska:sec:containers}
All graphs in this section are finite and undirected, may have loops, and have no multiple edges. A loop counts once in the number of neighbors of its vertex. An independent set contains neither both endpoints of an ordinary edge nor the vertex of a loop.

\begin{lemma}[Loop-aware regular-graph containers]\label{ska:lem:containers}
Let $G$ have $d$ vertices and exactly $k$ neighbors at every vertex, where $k\ge4$. Put
\begin{equation}\label{ska:eq:containerconstants}
t=\lfloor\sqrt{k}\rfloor,\qquad
\delta_k=\frac{t}{2(2k-t)}+\frac1{t+1},\qquad
b_k=\Hbin\left(\frac1{t+1}\right),
\end{equation}
where $\Hbin(p)=-p\log p-(1-p)\log(1-p)$. There is a family of at most $\exp(b_kd)$ vertex sets, each of size at most $d/2+\delta_kd$, such that every independent set is contained in one of them. Both $b_k$ and $\delta_k$ tend to zero as $k\to\infty$.
\end{lemma}
\begin{proof}
Fix a total order on vertices. Starting with the entire vertex set, repeatedly choose a vertex of maximum degree in the current induced graph, breaking ties by the fixed order. Stop when the induced maximum degree is at most $t$. To process an independent set $I$, if the queried vertex lies in $I$, record it in a fingerprint $F$ and remove its closed neighborhood; otherwise remove only that vertex. Every recorded vertex is loop-free, because it lies in $I$. Its induced degree is greater than $t$, so at least $t+1$ vertices are removed at this step. Therefore
\begin{equation}\label{ska:eq:fingerprint}
|F|\le d/(t+1).
\end{equation}
The final residual set $R$ has induced maximum degree at most $t$, and $I\subseteq F\cup R$.

Counting incidences directed from $R$ in the original regular graph gives
\[
k|R|\le t|R|+k(d-|R|).
\]
The internal incidences contribute at most $t|R|$; the crossing incidences contribute at most the total degrees outside $R$. A loop is internal and never crossing, so this estimate respects the loop convention. Thus
\[
|R|\le\frac{kd}{2k-t}
=\frac d2+\frac{t d}{2(2k-t)}.
\]
Combining with~\eqref{ska:eq:fingerprint} bounds $|F\cup R|$ as claimed.

The fingerprint reconstructs the full deterministic process: a pivot is treated as selected precisely when it is in $F$. Hence it determines $R$ and the container $F\cup R$. The number of containers is at most $\sum_{j\le d/(t+1)}\binom dj$. For $p=1/(t+1)\le1/2$ and $j\le pd$, one has
\[
p^j(1-p)^{d-j}\ge p^{pd}(1-p)^{(1-p)d}=e^{-d\Hbin(p)}.
\]
Summing the corresponding binomial probabilities proves
$\sum_{j\le pd}\binom dj\le e^{d\Hbin(p)}$. This proves the count and the lemma.
\end{proof}

\section{Dense windows and weighted preimages}\label{ska:sec:standardization}
Fix activities $(u,v,w)>0$ for letters $1,2,3$. In Sections~\ref{ska:sec:standardization}--\ref{ska:sec:criterion}, use
\begin{align}
\beta&=v+w,&\gamma&=u+v,&W&=u+v+w,\label{ska:eq:localnotation}\\
\Pi&=\beta(\beta+u),&T_0&=v^2+2uv,&
E&=\beta\gamma\sqrt{\frac{\beta+2u}{\beta+u}}.\label{ska:eq:pairfactors}
\end{align}
These are local weighted quantities, distinct from the polynomials $A,B,C,P$ in~\eqref{ska:eq:ABC}--\eqref{ska:eq:P}.

\subsection{The terminal reflection matching}
In any stressed word of length $L$, the terminal three forbids two ones at positions summing to $L$. Pair the positions $i,L-i$ for $1\le i<L$ and $i\ne L-i$. There are at most two exceptional positions: the terminal coordinate $L$, and the loop coordinate $L/2$ when $L$ is even. After removing these exceptions, all remaining coordinates lie in ordinary two-coordinate matching blocks. This matching remains a necessary constraint after any replacement that does not introduce a new one.

\subsection{Sparse threes}
Fix an integer $k\ge4$. Among words with fewer than $k^2$ threes, fix their positions and replace all threes by twos. The image has alphabet $\{1,2\}$ and retains the terminal matching constraint. Each ordinary pair has total weight $T_0=v^2+2uv$, while the at most two exceptions cost only a fixed factor. For fixed positions, replacing back has a weight ratio bounded in terms of $(u,v,w,k)$. There are at most $(L+1)^{k^2}$ choices of positions. Hence, with constants depending on the fixed activities and $k$,
\begin{equation}\label{ska:eq:sparse}
S_L^{\mathrm{sparse}}(u,v,w)
\le K_k(L+1)^{k^2}T_0^{L/2}.
\end{equation}
The harmless difference between $L/2$ and the number of ordinary pairs has been absorbed into $K_k$, since $T_0>0$ is fixed.

\subsection{A deterministic dense standardization}
Suppose there are at least $k^2$ threes. Let $n=\lceil L/k\rceil$ and partition $[L]=\{1,\ldots,L\}$ consecutively into intervals of length $n$, except possibly the last. There are at most $k$ intervals. Choose the rightmost interval containing at least $k$ threes; let $d$ be its rightmost three-position and let $Q$ be, for definiteness, its first $k$ three-positions. Then
\begin{equation}\label{ska:eq:Qwindow}
|Q|=k,\qquad Q\subseteq(d-n,d].
\end{equation}
Every later interval contains fewer than $k$ threes, and the chosen interval has none after $d$, so fewer than $k^2$ threes occur after $d$.

Standardize the word by replacing every positive position in $(d-n,d]$ by three, and every tail three at a position greater than $d$ by two. The image has a prefix of length $d$ on $\{1,2,3\}$ and a tail on $\{1,2\}$. No one was introduced, so the terminal reflection matching survives. The image need not itself be stressed; only the stated necessary constraints are used below. The number of overwritten positions is at most
\begin{equation}\label{ska:eq:overwritecount}
u_L=\lceil L/k\rceil+k^2.
\end{equation}
When $d<n$, the overwrite interval is intersected with positive positions, so the same bound remains valid.

Identify the prefix positions $1,\ldots,d$ with $\Z/d\Z$, where $d$ represents zero. Join $i$ to $q-i$ for every $q\in Q$. Each vertex has exactly $k$ neighbors; loops are allowed. The prefix one-set of the image is independent in this graph. To see this, two unchanged prefix ones both have index at most $d-n$. If their sum is congruent to $q\in Q$ modulo $d$, the wrapped possibility $q+d$ is excluded by
\[
q+d>2d-n>2d-2n.
\]
Their positive sum must therefore equal $q$, an original three-position, contrary to the original Kunz inequality. A changed prefix position is three and cannot be a one. This also covers repeated source indices and graph loops.

Lemma~\ref{ska:lem:containers} supplies a container $H$ for the prefix one-set, with
\begin{equation}\label{ska:eq:containeruse}
|H|\le d/2+\delta_kd,
\qquad\#\{H\}\le e^{b_kd}.
\end{equation}

\subsection{Every parameter and weight loss}
Before using any further relation, fix $d,Q$, and the tail overwrite set. This fixes the full overwrite set. At each overwritten position, the total possible original weight divided by the assigned image weight is at most
\begin{equation}\label{ska:eq:M}
M=\frac{u+v+w}{\min(v,w)}>1.
\end{equation}
Thus the weighted number of preimages of an image is at most $M^{u_L}$ times its image weight. This bound deliberately sums over all original letters, even those impossible in the fixed parameter class; such a relaxation is valid for an upper bound.

The choices of $d$, of its $k$ selected positions, and of fewer than $k^2$ tail positions are bounded by
\begin{equation}\label{ska:eq:parameterloss}
(L+1)^{k^2+k+1}.
\end{equation}
The full parameter and preimage loss is consequently at most
\begin{equation}\label{ska:eq:preimageloss}
K_k(L+1)^{k^2+k+1}\exp\left(\frac{\log M}{k}L\right).
\end{equation}
The extra factors $M^{k^2+1}$ have been included in $K_k$. In particular, after $k$ is fixed, all nonexponential losses are polynomials in $L$, and the exponential preimage rate tends to zero as $k\to\infty$.

\section{Matching blocks and the weighted criterion}\label{ska:sec:criterion}
Fix a standardized prefix length $d$ and a container $H$. A prefix coordinate is \emph{eligible} when it belongs to $H$. Relax the image family as follows: eligible prefix coordinates have alphabet $\{1,2,3\}$, ineligible prefix coordinates have alphabet $\{2,3\}$, and tail coordinates have alphabet $\{1,2\}$. In every ordinary terminal matching block, exclude $(1,1)$. Drop all other constraints, including fixed overwrite values. This is a product of independent blocks that contains every actual image under consideration.

Let $a$ be the number of prefix-prefix blocks, $b$ the number of crossing blocks, and $c$ the number of coordinates in tail-tail blocks. Let $e,f$ be the exceptional prefix and tail coordinate counts. Then
\begin{equation}\label{ska:eq:blockcounts}
2a+b+e=d,\qquad b+c+f=L-d,\qquad e+f\le2.
\end{equation}
The exact ordinary block weights are
\begin{center}
\begin{tabular}{@{}lll@{}}
\toprule
Block & Eligible prefix coordinates & Weight\\
\midrule
Prefix-prefix & $0$ & $\beta^2$\\
Prefix-prefix & $1$ & $\beta(\beta+u)$\\
Prefix-prefix & $2$ & $\beta(\beta+2u)$\\
Crossing & $0$ & $\beta\gamma$\\
Crossing & $1$ & $\beta\gamma+uv$\\
Tail-tail & none & $T_0$\\
\bottomrule
\end{tabular}
\end{center}
The exceptional coordinates contribute a factor bounded by a constant depending only on the activities. The successive eligibility gains are
\begin{equation}\label{ska:eq:gains}
L_1=\frac{\beta+u}{\beta},\qquad
L_2=\frac{\beta+2u}{\beta+u},\qquad
L_b=1+\frac{uv}{\beta\gamma},\qquad
L_1>L_2>L_b>1.
\end{equation}
For the only less immediate comparison, $L_2>L_b$ is equivalent to
\[
\beta\gamma>v(\beta+u),\qquad
\beta\gamma-v(\beta+u)=uw>0.
\]

\begin{lemma}[Two useful product bounds]\label{ska:lem:product}
The relaxed block product is bounded by each of
\begin{align}
K L_1^{\delta_kL}\Pi^a E^b T_0^{c/2},\label{ska:eq:sharpblocks}\\
K\Pi^a(\beta\gamma)^bT_0^{c/2}L_1^{|H|-a}.\label{ska:eq:crudeblocks}
\end{align}
Here $K$ depends only on the fixed activities, not on $k,L,d,H$.
\end{lemma}
\begin{proof}
The all-ineligible baseline is $\beta^{2a}(\beta\gamma)^bT_0^{c/2}$. By~\eqref{ska:eq:containeruse} and~\eqref{ska:eq:blockcounts}, the number of eligible prefix positions is at most
\[
a+b/2+e/2+\delta_kd.
\]
At the relaxed real budget $a+b/2$, grant the largest $a$ gains, one $L_1$ for each internal pair. Every remaining unit of eligibility costs at most $L_2$, including a crossing endpoint. Relaxing integral choices to fractional ones only increases the upper bound. Even if there are too few second internal endpoints to spend the remaining budget there, $L_2$ is still an upper bound on the smaller crossing gain. The resulting expression is
\[
\beta^{2a}(\beta\gamma)^bT_0^{c/2}L_1^aL_2^{b/2}
=\Pi^aE^bT_0^{c/2}.
\]
Increasing the budget by $e/2+\delta_kd$ costs at most $L_1^{e/2+\delta_kd}$. Since $e\le2$ and $d\le L$, this proves~\eqref{ska:eq:sharpblocks}. A smaller actual budget only reduces the product. Finally, bounding every eligibility gain by $L_1$ directly gives
\[
K\beta^{2a}(\beta\gamma)^bT_0^{c/2}L_1^{|H|},
\]
which equals~\eqref{ska:eq:crudeblocks}. Including exceptional eligible positions in $|H|$ is another harmless overestimate.
\end{proof}

\begin{proposition}[Positive-activity convergence criterion]\label{ska:prop:criterion}
For positive $(u,v,w)$, if
\begin{equation}\label{ska:eq:criterion}
\beta(\beta+u)<1,\qquad v^2+2uv<1,\qquad
\beta\gamma\sqrt{\frac{\beta+2u}{\beta+u}}<1,
\end{equation}
then $S_L(u,v,w)\le K\theta^L$ for some $K<\infty$ and $0<\theta<1$, uniformly in $L\ge1$.
\end{proposition}
\begin{proof}
Put $q_0=\max(\sqrt\Pi,\sqrt E,\sqrt{T_0})<1$. Because $2a+2b+c=L-e-f$,~\eqref{ska:eq:sharpblocks} is at most $K L_1^{\delta_kL}q_0^L$, with the at most two omitted coordinates absorbed into $K$. The square root of $E$ is correct: a crossing block has two coordinates. Summing containers and applying~\eqref{ska:eq:preimageloss} gives
\begin{equation}\label{ska:eq:densecriterion}
S_L^{\mathrm{dense}}(u,v,w)
\le K_k(L+1)^{k^2+k+1}
\left[q_0e^{b_k}L_1^{\delta_k}M^{1/k}\right]^L.
\end{equation}
The bracket tends to $q_0<1$ as $k\to\infty$. Choose one sufficiently large integer $k$ for which it is below one. This choice is made once, before $L$ tends to infinity. The fixed polynomial is absorbed into any slightly larger exponential base still below one. The sparse estimate~\eqref{ska:eq:sparse} has a strict gap because $T_0<1$, so the same conclusion holds after the two classes are added.
\end{proof}

\section{An exact certificate for the numerator}\label{ska:sec:numerator}
We now apply Proposition~\ref{ska:prop:criterion} after the activity transformation, at the rational point $z=33/50$. There is no floating-point inequality in this step.

\begin{proposition}[Exponential boundary summability]\label{ska:prop:numerator}
For some $K_N<\infty$ and $0<\theta_N<1$,
\begin{equation}\label{ska:eq:boundarydecay}
D_r(33/50)\le K_N\theta_N^r\qquad(r\ge1).
\end{equation}
Consequently $N$ is analytic on $|z|<33/50$ and the series in~\eqref{ska:eq:N} converges uniformly and absolutely on its closed disk.
\end{proposition}
\begin{proof}
At $z=33/50$, the three transformed activities in~\eqref{ska:eq:transform} are
\begin{align}
A^2&=\frac{8169809769}{15625000000}<\frac{523}{1000},\label{ska:eq:actcert1}\\
A^2C&=\frac{22377108957291}{39062500000000}<\frac{573}{1000},\label{ska:eq:actcert2}\\
A^2z^3&=\frac{293598453668553}{1953125000000000}<\frac{151}{1000}.\label{ska:eq:actcert3}
\end{align}
Use the dominating triple $(u,v,w)=(523,573,151)/1000$. Its exact pair factors are
\begin{equation}\label{ska:eq:triplecert}
\Pi=\frac{225707}{250000},\qquad
T_0=\frac{927687}{1000000},\qquad
E^2=\frac{108835743993}{121777343750}.
\end{equation}
The first two are below $19/20$, and $E^2<(19/20)^2$. All three inequalities in~\eqref{ska:eq:criterion} hold. The weighted polynomials $S_r$ have nonnegative coefficients, so coordinatewise domination and~\eqref{ska:eq:transform} prove~\eqref{ska:eq:boundarydecay}. For every $|z|\le33/50$, nonnegativity also gives $|D_r(z)|\le D_r(33/50)$. The resulting geometric majorant proves the asserted convergence and analyticity.
\end{proof}

At this point the separated class has a finite positive simple-pole amplitude. This does not yet settle the full stressed count: an uncontrolled complementary class could still have a divergent prefactor or another singularity at $\xi$. The next section proves a strict gap for exactly that complement.

\section{Early ones have an exponential critical gap}\label{ska:sec:early}
Fix $x=\xi$ and specialize the activities to $(x,x^2,x^3)$. In this section,
\begin{align}
\beta&=x^2+x^3,&\gamma&=x+x^2,&W&=x+x^2+x^3,\label{ska:eq:criticalletters}\\
\Pi&=\beta W=1,&T_0&=x^4+2x^3,&E&=\beta\gamma\sqrt{\frac{\beta+2x}{\beta+x}}.\label{ska:eq:criticalfactors}
\end{align}
The gains $L_1,L_2,L_b$ and $M$ retain their definitions from~\eqref{ska:eq:gains} and~\eqref{ska:eq:M}. Let
\begin{equation}\label{ska:eq:VL}
V_L=\sum_{\substack{w\text{ stressed},\ |w|=L\\
\exists q\le L/3:\ w_q=1}}x^{g(w)}.
\end{equation}

\begin{theorem}[Early-one suppression]\label{ska:thm:early}
There are $K_V<\infty$ and $0<\theta_V<1$ such that
\begin{equation}\label{ska:eq:earlygap}
V_L\le K_V\theta_V^L\qquad(L\ge1).
\end{equation}
\end{theorem}

\subsection{The strictly subcritical factors at the critical point}
Although $\Pi=1$, the other factors have strict margins:
\begin{equation}\label{ska:eq:criticalstrict}
\beta\gamma<1,\qquad T_0<1,\qquad E<1,
\qquad 1<L_b<L_2<L_1.
\end{equation}
The last chain was proved for all positive activities. Also $\gamma<W$, so $\beta\gamma<\beta W=1$. Since $x<2/3$, monotonicity gives $T_0<64/81<1$. To check $E$ without an unstated monotonicity assumption on a quotient, write
\[
E(x)^2=\frac{x^6(1+x)^4(x^2+x+2)}{x^2+x+1}.
\]
Its derivative is
\begin{equation}\label{ska:eq:Ederivative}
\frac{x^5(1+x)^3
(10x^5+26x^4+50x^3+51x^2+37x+12)}{(x^2+x+1)^2}>0.
\end{equation}
At $x=2/3$ the square equals $1120000/1121931<1$. This proves~\eqref{ska:eq:criticalstrict}.

\subsection{Reduction to fixed parameter classes}
Use the sparse/dense split and standardization of Section~\ref{ska:sec:standardization}, for a large fixed $k$ to be chosen. The sparse class already has a strict gap by~\eqref{ska:eq:sparse}. In the dense class, mark one original position $q\le L/3$ with $w_q=1$. Summing over all possible marks introduces at most a factor $L$ and is used only for an upper bound.

Fix $q,d,Q$, the full overwrite set $F$, and a container $H$. These data are fixed before any product probability is considered. The image one-set lies in $H$, and the matching-block relaxation of Section~\ref{ska:sec:criterion} contains every actual image. The two product bounds are now
\begin{align}
K L_1^{\delta_kL} E^b T_0^{c/2},\label{ska:eq:criticalsharp}\\
K(\beta\gamma)^bT_0^{c/2}L_1^{|H|-a}.\label{ska:eq:criticalcrude}
\end{align}
The parameters still satisfy~\eqref{ska:eq:blockcounts}, and $|F|\le u_L$ from~\eqref{ska:eq:overwritecount}. Set $\delta=1/48$. The following three cases are exhaustive.

\subsection{Case one a short standardized prefix}
Suppose $d\le(1-\delta)L$. Put $\sigma=\max(E,\sqrt{T_0})<1$. By~\eqref{ska:eq:blockcounts},
\[
b+c=L-d-f\ge\delta L-2.
\]
Consequently~\eqref{ska:eq:criticalsharp} is at most
\begin{equation}\label{ska:eq:caseone}
K L_1^{\delta_kL}\sigma^{\delta L-2}.
\end{equation}
Before restoring the vanishing container and preimage rates, this has the fixed positive exponential saving $-\delta\log\sigma$.

\subsection{Case two a long prefix and small container}
Suppose $d>(1-\delta)L$ and $|H|<5L/12$. Since $b\le L-d$ and $e\le2$,
\[
a=\frac{d-b-e}{2}\ge\frac{2d-L-2}{2}
>\left(\frac12-\delta\right)L-1.
\]
It follows that $a-|H|>L/16-1$. The factors $\beta\gamma,T_0$ are below one, so~\eqref{ska:eq:criticalcrude} is at most
\begin{equation}\label{ska:eq:casetwo}
K L_1^{-L/16+1}.
\end{equation}
This gives the fixed positive saving $(\log L_1)/16$.

\subsection{Case three a long prefix and large container}
Suppose $d>(1-\delta)L$ and $|H|\ge5L/12$. At most $q$ members of $H$ lie above $d-q$, so there are at least
\[
|H|-q\ge L/12
\]
eligible source positions $i\in H$ satisfying $i+q\le d$. Discard a source if either endpoint $i,i+q$ belongs to $F$, if either is an exceptional matching coordinate, or if the two endpoints belong to the same ordinary matching block. The first exclusion costs at most $2|F|\le2u_L$, the second at most four, and the last at most one, since it requires $2i+q=L$. Thus the number $m$ of retained relations satisfies
\begin{equation}\label{ska:eq:mrelations}
m\ge L/12-2u_L-5.
\end{equation}

Every actual image $v$ in the fixed parameter class obeys
\begin{equation}\label{ska:eq:translated}
(v_i,v_{i+q})\ne(1,3)
\end{equation}
for each retained source. Indeed, the original marked value was $w_q=1$, and the original inequality gives $1+w_i\ge w_{i+q}$. The queried endpoints were unchanged, so the same forbidden pair holds for their image values. This argument uses the \emph{original} value at $q$, not the image value $v_q$. The marked position itself may have been overwritten. Nor is any assumption made that every eligible source actually carries one.

Now form a multigraph whose vertices are the ordinary terminal matching blocks and whose edges are the retained directed relations $i\to i+q$, regarded as undirected edges for selection. There are no loops, since same-block relations were discarded. A coordinate is the source of at most one relation and the target of at most one. Each block has two coordinates, so its degree, counting multiplicities, is at most four. Repeatedly select the first remaining edge in a fixed order and delete all edges meeting either endpoint block. Each selection deletes at most eight edges, and hence selects at least $m/8$ pairwise block-disjoint relations.

This selection is determined entirely by $q,H,F,d,L$ and the fixed order. It is completed \emph{before} drawing any letters from the product model. Parallel relation-edges cause no difficulty.

Normalize the relaxed product of ordinary block weights to a probability measure. For a selected relation, its source is eligible, so the block assignment
\[
(\text{source}=1,\ \text{partner}=2)
\]
is allowed and has weight $x^3$. Its target is a prefix coordinate, so the assignment
\[
(\text{target}=3,\ \text{partner}=2)
\]
is allowed and has weight $x^5$, irrespective of target eligibility or of the partner's side. Each block's total weight is at most $W^2$, because every allowed pair is in $\{1,2,3\}^2$. The source and target blocks are distinct. Therefore
\begin{equation}\label{ska:eq:eta}
\Pr(v_i=1,\ v_{i+q}=3)\ge\eta:=\frac{x^8}{W^4}\in(0,1).
\end{equation}
Distinct selected relations involve disjoint blocks, so these violation events are independent in the relaxed product measure. Their exclusion multiplies the product bound by at most
\begin{align}
(1-\eta)^{m/8}
&\le K_k\exp\left[-\frac\eta8\left(\frac1{12}-\frac2k\right)L\right]\label{ska:eq:penalty}\\
&\le K_k e^{-\eta L/128}\qquad(k\ge96).\nonumber
\end{align}
The ceiling error and the fixed $k^2$ term in $u_L$ are absorbed into $K_k$. When the lower bound in~\eqref{ska:eq:mrelations} is negative, the same formula remains a harmless loose bound after this constant is included. Together with~\eqref{ska:eq:criticalsharp} and $E,T_0<1$, the case-three image weight is at most
\begin{equation}\label{ska:eq:casethree}
K_kL_1^{\delta_kL}e^{-\eta L/128}.
\end{equation}

The independence used here belongs only to the enlarged product of matching blocks. We never assert independence in the actual Kunz ensemble, and never condition the product measure on Kunz validity, overwrite values, or the value at the original mark. These are precisely the constraints that can be dropped for an upper bound before imposing the necessary exclusions~\eqref{ska:eq:translated}.

\subsection{One parameter choice for all cases}
Define the positive fixed saving
\begin{equation}\label{ska:eq:tau}
\tau=\min\left(-\frac{\log\sigma}{48},
\frac{\log L_1}{16},\frac\eta{128}\right)>0.
\end{equation}
The combined container, eligibility, and preimage exponential rate is at most
\begin{equation}\label{ska:eq:epsilonk}
\eps_k=b_k+\delta_k\log L_1+\frac{\log M}{k}\longrightarrow0.
\end{equation}
Choose one integer $k\ge96$ with $\eps_k<\tau/4$. This choice works for all three cases. After multiplying their image estimates by all parameter and preimage bounds, and summing containers and the at most $L$ marks, the dense contribution is at most
\[
K_k(L+1)^{k^2+k+2}e^{-3\tau L/4}
\le K'_k e^{-\tau L/2}.
\]
The second inequality follows because a fixed polynomial times $e^{-\tau L/4}$ is bounded over positive integers. The sparse contribution~\eqref{ska:eq:sparse} is at most another constant times $T_0^{L/4}$. Thus~\eqref{ska:eq:earlygap} holds, for example with
\[
\theta_V=\max(e^{-\tau/2},T_0^{1/4})<1
\]
and a sufficiently large finite $K_V$. This completes the proof of Theorem~\ref{ska:thm:early}.

The ordering is important: fix the activities and all positive margins, choose a single sufficiently large $k$, and only then let $L$ increase. The proof makes no exchange of a growing standardization parameter with the word length.

\section{The analytic disk and the positive simple pole}\label{ska:sec:analytic}
We now prove Theorem~\ref{ska:thm:main}. Let $\cV(z)$ be the genus generating function of the early-one class in~\eqref{ska:eq:VL}. Every length-$L$ word has genus at most $3L$. For $y>x=\xi$,
\[
\sum_{\substack{w\in V\,;\ |w|=L}}y^{g(w)}
\le (y/\xi)^{3L}V_L
\le K_V[\theta_V(y/\xi)^3]^L.
\]
For example, $y=\xi\theta_V^{-1/6}>\xi$ makes the bracket $\sqrt{\theta_V}<1$. Hence $\cV(y)<\infty$, and its nonnegative coefficient series is analytic for $|z|<y$.

Proposition~\ref{ska:prop:numerator} makes $N$ analytic for $|z|<33/50$. The exact formal decomposition~\eqref{ska:eq:proofmap}, initially valid near zero, therefore gives a meromorphic continuation of $\cS$ to the intersection of these disks.

There is no zero of $1-P(z)$ with $|z|<\xi$, since $|P(z)|\le P(|z|)<1$. On $|z|=\xi$, the equation $P(z)=1$ implies equality in the triangle inequality, and each nonzero monomial of $P(z)$ must have positive real phase. Writing $z=\xi e^{i\vartheta}$ gives $e^{3i\vartheta}=e^{4i\vartheta}=1$, and thus $e^{i\vartheta}=1$. The only boundary zero is $z=\xi$. It is simple, because $P'(\xi)>0$.

Since the polynomial has finitely many roots, one can choose $R_s>\xi$ smaller than $33/50$, smaller than $y$, and smaller than the modulus of every other zero of $1-P$. Near $\xi$,
\[
1-P(z)=\xi P'(\xi)(1-z/\xi)+O((z-\xi)^2).
\]
Therefore the function
\begin{equation}\label{ska:eq:analyticremainder}
H(z)=\frac{N(z)}{1-P(z)}+\cV(z)
-\frac{N(\xi)}{\xi P'(\xi)}\frac1{1-z/\xi}
\end{equation}
has a removable singularity at $\xi$ and is analytic on $|z|<R_s$. For any fixed $R'$ with $\xi<R'<R_s$, Cauchy's estimate yields
\[
|[z^g]H(z)|\le \max_{|z|=R'}|H(z)|\,(R')^{-g}.
\]
Set $\rho_2=(R')^{-1}<\rho$ to obtain~\eqref{ska:eq:main}. Positivity and finiteness follow from the convergent nonnegative numerator and its strictly positive baseline:
\begin{equation}\label{ska:eq:positivity}
0<\frac{\xi^3(1+A(\xi))}{\xi P'(\xi)}\le C_s<\infty.
\end{equation}
No Tauberian argument, ratio-limit hypothesis, or extrapolation of coefficients is involved.

\begin{corollary}[Consequences for the stressed sequence]\label{ska:cor:stress}
For some $0<\lambda<1$,
\begin{equation}\label{ska:eq:stressconsequences}
\frac{s_g}{\rho^g}=C_s+O(\lambda^g),\qquad
\frac{s_{g+1}}{s_g}=\rho\bigl(1+O(\lambda^g)\bigr),\qquad
s_g=g^{o(1)}\rho^g.
\end{equation}
The sequence is positive and strictly increasing for all sufficiently large $g$, and its $g$th root tends to $\rho$.
\end{corollary}
\begin{proof}
Take $\lambda=\rho_2/\rho$. Divide~\eqref{ska:eq:main} by $C_s\rho^g$, whose constant is positive. The ratio conclusion follows by division of $1+O(\lambda^{g+1})$ by $1+O(\lambda^g)$. Taking logarithms gives the root and $g^{o(1)}$ conclusions. Finally,
\[
s_{g+1}-s_g=C_s(\rho-1)\rho^g+O(\rho_2^g)>0
\]
eventually, because $\rho>1$.
\end{proof}

\section{Two poles for the depth at most three count}\label{ska:sec:twopole}
Retain $t_g,h_g,n_g$ and the exact filter of Section~\ref{ska:sec:kunz}. Write
\begin{equation}\label{ska:eq:Slead}
S_{\mathrm{lead}}=\frac{\varphi}{\sqrt5}
\left(1+\cS(\varphi^{-1})\right)
=\frac{\varphi}{\sqrt5}\left(1+\sum_{g\ge0}s_g\varphi^{-g}\right),
\qquad
\kappa=\frac1{\xi+\xi^2-1}.
\end{equation}
The series converges because $\varphi^{-1}<\xi$, and $\kappa>0$. This is the usual leading constant from the Fibonacci filter, with a subscript only to distinguish it from the generating function $\cS$.

\begin{corollary}[The shallow residual]\label{ska:cor:twopole}
For some $0<\rho_3<\rho$,
\begin{equation}\label{ska:eq:twopole}
t_g=S_{\mathrm{lead}}\varphi^g-\kappa C_s\rho^g+O(\rho_3^g).
\end{equation}
In particular, $(t_g-S_{\mathrm{lead}}\varphi^g)/\rho^g\to-\kappa C_s<0$. For all semigroups the conclusion is exactly
\begin{equation}\label{ska:eq:total}
n_g=S_{\mathrm{lead}}\varphi^g+h_g-\kappa C_s\rho^g+O(\rho_3^g).
\end{equation}
\end{corollary}
\begin{proof}
The roots of $1-z-z^2$ are $r=\varphi^{-1}$ and $-\varphi$. At $r$, the coefficient of $(1-z/r)^{-1}$ in the filter is
\[
\frac{1+\cS(r)}{r(1+2r)}
=\frac{\varphi}{\sqrt5}(1+\cS(r))=S_{\mathrm{lead}}.
\]
At $\xi$, the filter denominator is nonzero. The coefficient of $(1-z/\xi)^{-1}$ is
\[
\frac{C_s}{1-\xi-\xi^2}=-\kappa C_s.
\]
Choose $R$ with $\xi<R<\min(R_s,\varphi)$. After subtracting these two simple-pole terms, the function is analytic for $|z|<R$. The negative Fibonacci pole lies outside this disk. Cauchy's estimate on a slightly smaller radius greater than $\xi$ proves~\eqref{ska:eq:twopole}, and $n_g=t_g+h_g$ proves~\eqref{ska:eq:total}.
\end{proof}

The explicit $h_g$ in~\eqref{ska:eq:total} is essential. This article does not show that $h_g/\rho^g$ converges, determine $h_g/s_g$, or prove a sign or noncancellation for the total residual. If a finite limit $C_h$ were established separately, its total secondary coefficient would be $C_h-\kappa C_s$; equality would require a further analysis on a smaller scale.

\section{An inverse threshold and fixed-order corrections}\label{ska:sec:inverse}
The finite amplitude improves inversion beyond a statement of exponential root growth. Fix constants from~\eqref{ska:eq:main} and put
\begin{equation}\label{ska:eq:inverseparams}
\lambda=\rho_2/\rho\in(0,1),\qquad
\delta_*=-\frac{\log\lambda}{\log\rho}>0,\qquad
x_0(y)=\frac{\log(y/C_s)}{\log\rho}.
\end{equation}
For sufficiently large positive real $y$, define the first threshold
\[
G(y)=\min\{g\ge0:s_g\ge y\}.
\]
It exists because $s_g\to\infty$. Enlarge the range of $y$ so that it exceeds all counts in the finite initial segment before eventual strict monotonicity.

\begin{corollary}[A shrinking integer-boundary bracket]\label{ska:cor:inverse}
There are finite constants $K_*,Y$ such that, for all real $y\ge Y$,
\begin{equation}\label{ska:eq:ceilingbracket}
\left\lceil x_0(y)-K_*y^{-\delta_*}\right\rceil
\le G(y)\le
\left\lceil x_0(y)+K_*y^{-\delta_*}\right\rceil.
\end{equation}
If $\dist(x_0(y),\Z)>K_*y^{-\delta_*}$, the two endpoints coincide and $G(y)=\lceil x_0(y)\rceil$. If the input is promised to equal $s_g$ for a sufficiently large integer $g$, then $g$ is the unique nearest integer to $x_0(y)$.
\end{corollary}
\begin{proof}
Write $s_g=C_s\rho^g(1+e_g)$ with $|e_g|\le M_0\lambda^g$. Eventually $|e_g|\le1/2$, so
\begin{equation}\label{ska:eq:logerror}
\left|\frac{\log s_g}{\log\rho}-g-\frac{\log C_s}{\log\rho}\right|
\le\frac{2M_0}{\log\rho}\lambda^g.
\end{equation}
For integers satisfying $|g-x_0(y)|\le2$, the right side is bounded by $K_0y^{-\delta_*}$ with a fixed $K_0$, because
\[
\lambda^g\le\lambda^{-2}\lambda^{x_0(y)}
=\lambda^{-2}C_s^{\delta_*}y^{-\delta_*}.
\]
Choose $K_*>K_0$. Increase $Y$ until $\eps=K_*y^{-\delta_*}<1/4$ and all integers within distance two of $x_0(y)$ lie in the increasing tail. Put
\[
g_- =\lceil x_0(y)-\eps\rceil-1,
\qquad g_+=\lceil x_0(y)+\eps\rceil.
\]
Both are within distance two of $x_0(y)$, and
$g_-<x_0(y)-\eps$ while $g_+\ge x_0(y)+\eps$.
The strict slack $K_*>K_0$ in~\eqref{ska:eq:logerror} gives $s_{g_-}<y<s_{g_+}$, including when either shifted endpoint is itself an integer. Eventual monotonicity, together with exclusion of the finite initial segment, yields $g_-+1\le G(y)\le g_+$, which is~\eqref{ska:eq:ceilingbracket}.

For the exact-count promise $y=s_g$,
\[
x_0(s_g)=g+\frac{\log(1+e_g)}{\log\rho}
=g+O(\lambda^g)=g+O(y^{-\delta_*}).
\]
Eventually the error is strictly less than one half, proving unique nearest-integer recovery. This last conclusion uses the exact-count promise; nearest-integer rounding is generally incorrect for arbitrary thresholds.
\end{proof}

Near an integer boundary, the bracket~\eqref{ska:eq:ceilingbracket} must be retained. The error estimate does not choose a side of that boundary. The constants in this inversion result are existence quantities; without a certified evaluation of $C_s$ and the remainder, it is not a practical numerical inverse algorithm.

\begin{corollary}[Every fixed inverse-power correction is zero]\label{ska:cor:allorders}
For every fixed integer $N\ge1$,
\begin{equation}\label{ska:eq:allorders}
s_g=C_s\rho^g\bigl(1+o(g^{-N})\bigr).
\end{equation}
Consequently, in any fixed-order expansion
\[
s_g=C_s\rho^g\left(1+\frac{c_1}{g}+\cdots+
\frac{c_N}{g^N}+o(g^{-N})\right),
\]
all coefficients $c_1,\ldots,c_N$ vanish.
\end{corollary}
\begin{proof}
The relative error in~\eqref{ska:eq:main} is $O(\lambda^g)$, which is $o(g^{-N})$ for every fixed $N$. If a displayed coefficient were nonzero, multiply by the smallest corresponding power $g^j$ after subtracting one and take the limit. The exponential remainder tends to zero, contradicting that nonzero coefficient.
\end{proof}

This is an all-fixed-algebraic-orders statement, not a full subdominant transseries. It does not locate the next singularity, determine its amplitude, rule out oscillations at smaller exponential scales, or identify the optimal remainder base.

\begin{remark}[{[write]} The inverse and the transseries volume, 7~October 2026]\label{ska:rem:transseries}
Compared statement by statement with the repository's transseries
volume~\cite{TSvol}; the volume's symbols are not used here.
\begin{enumerate}\raggedright
\item \emph{The growth.} If Theorem~\ref{ska:thm:main} holds, then along the
integers $\log s_g=\log C_s+g\log\rho+O(\lambda^g)$: the
exponential--power model of \file{p0:def:model} with exponent one, no
logarithmic term and zero correction series, every correction being smaller
than every power of $g$ (which is Corollary~\ref{ska:cor:allorders}). This is
an exact instance of the model identity along the sequence, conditional on
the theorem.
\item \emph{The inverse centre.} $x_0(y)$ of~\eqref{ska:eq:inverseparams} is
the exact solution of the volume's dominant block \file{p0:thm:lambert-core}
in its degenerate case without logarithmic term (part~(1)), with slope
$\log\rho$ and target $\log(y/C_s)$: an exact instance.
\item \emph{The thresholds.} $G(y)$ is the volume's staircase
(\file{p0:def:three-inverses}(1)) of the eventually increasing~$s_g$. The model
$C_s\rho^x$ does not take the values $s_g$ at the integers, so it is not an
admissible interpolation: the bracket~\eqref{ska:eq:ceilingbracket}, with its
clause for $\dist(x_0,\Z)$, is an analogue of \file{p0:thm:staircase}(2), and
the nearest-integer recovery of a promised exact count an analogue of
\file{p0:thm:staircase}(3), not instances. Both are proved directly, and both
inherit the conditional status of Theorem~\ref{ska:thm:main}.
\end{enumerate}
\end{remark}

\section{Bounded exact checks and reproducibility}\label{ska:sec:checks}
The mathematical proof above is independent of computation. The accompanying original Python companion checks finite instances and exact algebra to detect implementation or transcription errors. All decisive arithmetic is integer or rational arithmetic. No numerical estimate of $C_s$ is supplied.

\subsection{What is checked}
The companion independently constructs stressed words, boundary polynomials, and the length-refined renewal coefficients. It compares the activity transformation both as a genus-polynomial identity and at unrelated positive rational activity triples. It also checks the local eligibility optimization, loop-aware fingerprints, dense standardization, unchanged-endpoint translated exclusions, the deterministic block matching, and the exact lower bound for forbidden block marginals. The certificate at $33/50$ and the critical rational comparisons are checked exactly.

The default run checks boundary polynomial identities through $r=10$, the full length-and-genus renewal through $L=12$, and $30$ general-activity identities. It checks $7{,}056$ eligibility allocations and $1{,}470$ exact optimum budgets; $502$ cyclic graphs, including $476$ with loops; and $4{,}330$ independent-set/container instances. The dense construction is exercised on $41{,}878$ images, with $11{,}358$ original early marks, including $872$ overwritten marks. It verifies $9{,}123$ retained translated relations and $7{,}768$ selected block-disjoint relations, as well as $54$ local marginal comparisons and three explicit four-block product checks. These figures refer to the fixed default caps documented in the README, rather than an unbounded test claim.

The tests include endpoints and repeated source indices, empty initial genus coefficients, both parity cases of the terminal matching, and small structural examples in which an original marked coordinate is overwritten. Small standardization parameters are used only to exercise constructions in bounded finite tests; the theorem chooses a much larger fixed parameter to obtain its asymptotic savings.

\subsection{Caps and the trust boundary}
Public computation parameters are strictly bounded before enumeration. Validation uses explicit runtime checks that remain active under \texttt{python -O}; it does not use removable assertions as an input-safety gate. Textual integer input is subject to a 640-digit cap before conversion, in addition to mathematical range caps. This digit cap is a parsing-resource limit, not a mathematical precision claim and not a license to run an enumeration of a size specified by a long integer. The builder and tests never install packages or access the network automatically.

The offline builder checks an exact public-file allowlist, snapshots its sources, verifies content consistency, and writes outside the source directory. Its Python-managed file handling checks paths, symbolic links, hard links, and relevant file identities. Those controls are not a general sandbox for external executables. In particular, the installed TeX engine is explicitly outside the Python race-safety boundary: compilation executes the trusted supplied TeX with shell escape disabled, and the bundle does not claim to sandbox a malicious TeX program or a hostile concurrent filesystem adversary during external compilation.

The reproducible package contains the editable article, its PDF, the companion and tests, a short source map, and the builder documentation. It does not redistribute third-party papers, prior supporting reports, private working records, or audit notes. The complete mathematical argument needed for this theorem is in this article. A manifest pins the public members, and a separate checksum pins the actual ZIP bytes. Normal and optimized-Python replays check that validation survives optimization and that the declared outputs agree. Exact commands and tested resource caps are specified in the package README.

Finite tests validate only their stated bounded domains. They do not prove the universal weighted criterion, justify an asymptotic extrapolation, or replace the three-case early-one argument.
\begin{writenote}[The shipped files, 7 October 2026]
In ProveIt the PDF and the manifest are not shipped and the delivered README
is replaced; the companion, the tests and the builder are
\file{code/companion-exact_checks.py}, \file{code/tests-test_companion.py},
\file{code/tests-test_build.py} and \file{code/build.py}. The builder checks
the exact delivered seven-file layout (\file{companion/}, \file{tests/}), and
the companion test imports the companion by its delivered path, so both run
only in a re-extracted archive; the report README gives commands.
\end{writenote}

\begin{remark}[{[write]} The OEIS, 7~October 2026]\label{ska:rem:oeis}
The source cites no OEIS entry. The count $n_g$ of all numerical semigroups of
genus~$g$ is OEIS A007323, revision \#263 (18 August 2026), by Don Zagier
(1994), which links Zhu's paper and displays Bras-Amor\'os's conjectures,
among them $a(n)\ge a(n-1)+a(n-2)$ (Zhu's Conjecture~1.2); the total
formula~\eqref{ska:eq:total} does not decide it, as the paragraph after
Corollary~\ref{ska:cor:twopole} says. Searches of 7~October 2026 for initial
segments of $s_g$ ($1,0,1,3,2,4,9,12,20,32,\ldots$ from $g=3$) and of $t_g$
($1,1,2,4,6,11,20,33,57,99,\ldots$) found no entry; neither sequence is in the
OEIS. No OEIS edit is part of this work.
\end{remark}

\section{Further questions}\label{ska:sec:questions}
\subsection{Depth four and the full secondary residual}
The principal remaining enumeration question is whether $h_g=C_h\rho^g+o(\rho^g)$ for some finite positive $C_h$, or whether a genuinely different prefactor occurs. The early-one argument exploits that every unwrapped three-letter violation is $(1,1,3)$. Depth four introduces additional forbidden patterns and a wrapped constraint, so the present block penalty does not transfer without a new structural analysis. A useful next goal is an exact depth-four decomposition whose boundary series has a checkable convergence criterion. Even if $C_h$ exists, proving $C_h\ne\kappa C_s$ is a separate question.

\subsection{A certified and practical amplitude calculation}
The positive truncations
\[
C_s^{(R)}=\frac{\xi^3(1+A(\xi))+\sum_{r=1}^{R}D_r(\xi)}
{\xi P'(\xi)}
\]
increase to $C_s$. They are rigorous lower approximants if the algebraic evaluation of $\xi$ is itself enclosed rigorously. A useful numerical method also needs a computable upper tail bound. The existence proof provides exponential summability but potentially enormous constants from the standardization parameter and polynomial absorption. Improving this tail control, perhaps by combining a finite transfer computation with a quantitative weighted contraction, is needed before moderate truncations can certify many digits of the amplitude.

\subsection{Explaining the finite-genus regime}
The exponential asymptotic error does not by itself explain the apparent positive slope of finite-genus ratio plots. Determining when the early-one and boundary tails become negligible, and which families dominate before that point, would connect the theorem to observable finite data. This requires estimates with controlled constants, rather than fitting an eventual exponent to a bounded initial range.

\subsection{Subdominant singularities and optimal error rates}
The argument isolates one simple pole and proves an analytic neighborhood beyond it. The next obstruction might arise from another zero of $1-P$, from the boundary numerator, or from the early-one series. Locating it would determine whether smaller terms are pure exponentials, oscillatory contributions, or more complicated singular terms. The arbitrary admissible bases $\rho_2$ and $\rho_3$ used here are not asserted to be optimal.

\subsection{General activity domains}
The inequalities in Proposition~\ref{ska:prop:criterion} are sufficient conditions, not a characterization of all convergent activity triples. The invariant $\gamma/\alpha$ and the exact transformation~\eqref{ska:eq:Tmap} suggest studying the true convergence boundary of $\sum_rS_r(u,v,w)$. A sharper description could strengthen amplitude tail bounds and distinguish geometric mechanisms that are hidden by a uniform container relaxation. Any such use must retain the distinction between the exact separated-class renewal and the full stressed series.
\begin{writenote}[Added 7 October 2026, under Vladimir's standing rule of 4 October 2026]
These five directions are the source's own and stay open. The write adds two.
\emph{Independent verification} of Theorem~\ref{ska:thm:main}, which the source
requests: the write's reading found no error, but no second reader has
checked the all-length argument, and Zhu's finite data are far from the
asserted limit (note at the end of Section~1.2). \emph{The boundary series
at the critical point.} With the transformed activities
$(u,v,w)=(A^2,A^2C,A^2\xi^3)(\xi)=(0.5228\ldots,0.5727\ldots,0.1502\ldots)$,
whose factors $\Pi=0.9008\ldots$, $T_0=0.9269\ldots$, $E=0.9438\ldots$ satisfy
the criterion with the rate bound $\max(\sqrt\Pi,\sqrt E,\sqrt{T_0})=0.9715\ldots$,
the write's floating values $D_r(\xi)=0.4386$, $0.4707$, $0.4989$, $0.5353$, $0.5664$, $0.5907$ for $r=23,25,\ldots,33$ and $0.3003$, $0.3277$, $0.3472$, $0.3665$, $0.3886$ for $r=24,26,\ldots,32$ are still increasing (by factors
between $1.04$ and $1.10$ every two steps for $21\le r\le33$, slowly
decreasing), and the lower approximants $C_s^{(R)}$ are $0.4490$, $1.0798$, $2.0415$, $2.4159$ for $R=10,20,30,33$ (the baseline term alone gives $0.1199$). The proof's polynomial factors allow such a transient; locating
the turning point, or a certified tail bound, would test the theorem
quantitatively. No claim of the source was found false.
\end{writenote}

\clearpage
\begin{thebibliography}{9}
\bibitem{Bacher}
Roland Bacher, \emph{Generic numerical semigroups}, preprint, 2021.
\href{https://arxiv.org/abs/2105.04200}{arXiv:2105.04200}.
Transition-region discussion: \href{https://arxiv.org/html/2105.04200v1\#S15.SS4}{Section~15.4}.

\bibitem{Report270}
\emph{Sharp Secondary Exponential Growth of Numerical Semigroups}, Research report 270, 6 October 2026. Context for the dense-window and matching framework; no theorem from this report is required as a black-box input here.

\bibitem{Report271}
\emph{Consecutive Ratios at the Secondary Scale of Numerical Semigroup Enumeration}, Research report 271, 6 October 2026. Context for terminal-matching optimization and the separated-boundary investigation; not a proof dependency of the present amplitude theorem.

\bibitem{Samotij}
Wojciech Samotij, \emph{Counting independent sets in graphs}, European Journal of Combinatorics \textbf{48} (2015), 5--18.
\href{https://doi.org/10.1016/j.ejc.2015.02.005}{doi:10.1016/j.ejc.2015.02.005}.
\href{https://arxiv.org/abs/1412.0940}{arXiv:1412.0940}.

\bibitem{Zhu}
Daniel G. Zhu, \emph{Sub-Fibonacci behavior in numerical semigroup enumeration}, Combinatorial Theory \textbf{3}(2), article 10, 2023.
\href{https://doi.org/10.5070/C63261988}{doi:10.5070/C63261988}.
\href{https://arxiv.org/html/2202.05755v3}{arXiv version 3, 7 September 2023}.
Stressed words and the filter: Section~3; dense-window standardization: Sections~5.4--5.5; polynomial-prefactor formulation: Section~7.3(c); secondary cancellation discussion: Section~7.4.
\bibitem{TSvol}
ProveIt, \emph{Transseries and inversion}, \path{Analysis/Transseries/docs/series-and-transseries/Transseries_And_Inversion/transseries_and_inversion.tex}; Part~0 (\file{p0:def:model}, \file{p0:thm:lambert-core}, \file{p0:def:three-inverses}, \file{p0:thm:staircase}). [Added in the write.]
\end{thebibliography}
\end{document}
